\section{LEVEL 7 · Eindmissie}

\begin{missie}
Missie: combineer alles, werk als echte onderzoeker en leg je redenering uit. (LPD 1 – 4)
\end{missie}

\begin{doel}
\textbf{DOEL} \quad Na dit level kan je:

☐ een \textbf{fenomeen uit de realiteit beschrijven} met wiskundige begrippen (LPD 1);

☐ een probleem \textbf{aanpakken met een heuristiek}, wiskundig vertalen en terugvertalen met een antwoordzin (LPD 2);

☐ \textbf{ICT} gebruiken voor berekeningen en grafieken (LPD 3);

☐ je \textbf{redenering} schriftelijk en mondeling \textbf{uitleggen} en fouten in redeneringen verbeteren (LPD 4).
\end{doel}

\subsection*{Theorie 7.1 · De aanpak van een probleem}

\begin{theorie}
Kom je er niet meteen uit? Gebruik een heuristiek:

\begin{enumerate}
\item \textbf{Gegeven en gevraagd} noteren: wat weet ik, wat zoek ik?
\item \textbf{Schets of tekening} maken (een Gausskromme, een tabel, een tijdlijn).
\item Het probleem \textbf{opdelen in deelproblemen}.
\item \textbf{Wiskundig vertalen} (mathematiseren): kies het juiste begrip en de juiste formule of ICT-functie.
\item \textbf{Terugvertalen} (demathematiseren): schrijf een antwoordzin in de taal van de context.
\item \textbf{Controleren}: is mijn antwoord realistisch?
\end{enumerate}
\end{theorie}

\subsection*{Eindmissie 1 · Mini-onderzoek: hoeveel slapen jongeren?}

De leerlingenraad van een school met 600 leerlingen in de derde graad wil weten hoeveel uur leerlingen slapen. Ze bevragen 30 leerlingen. Hun slaapduur (in uur) op een schoolnacht:

6,1 – 7,6 – 7,5 – 7,3 – 5,1 – 8,1 – 6,0 – 6,7 – 7,1 – 7,3 – 7,7 – 8,0 – 7,1 – 7,4 – 6,9 – 7,2 – 7,9 – 8,5 – 6,4 – 7,5 – 6,8 – 6,9 – 7,7 – 6,4 – 6,9 – 8,3 – 5,8 – 5,6 – 9,3 – 8,5

\textbf{a)} \emph{(LPD 1)} Noteer de populatie, de steekproef, de variabele en het soort variabele. Wat zijn $N$ en $n$? \lijn

\textbf{b)} \emph{(LPD 3)} Bereken met ICT $\overline{x}$, $s$, de mediaan, $Q_1$ en $Q_3$. Maak met ICT een histogram (klassenbreedte 0,5 uur). \lijn

\textbf{c)} \emph{(LPD 13)} Is de normale verdeling een goed model? Verklaar met je histogram en met $\overline{x}$ en de mediaan. Noteer het model $N(\mu,\sigma)$. \lijn

\textbf{d)} \emph{(LPD 14)} Bereken de kans dat een leerling minder dan 6 uur slaapt. Hoeveel van de 600 leerlingen verwacht je dat? Schrijf een antwoordzin. \lijn

\textbf{e)} \emph{(LPD 11)} De 30 leerlingen zijn de leerlingen die in de pauze in de leefruimte zaten. Is de steekproef representatief? Hoe kan het beter? \lijn

\textbf{f)} \emph{(LPD 12)} Een leerling zegt: ‘Wie meer slaapt, haalt betere punten, dus meer slapen zorgt voor betere punten.’ Wat is hier mis? Bedenk twee alternatieve verklaringen en een manier om causaliteit te onderzoeken. \lijn

\textbf{g)} \emph{(LPD 4)} Schrijf voor de directie een besluit van vier à vijf zinnen waarin je de resultaten, de beperkingen van de steekproef en een advies verwerkt. \lijn

\begin{oplossing}
\textbf{Oplossing Eindmissie 1}

a) Populatie: de 600 leerlingen van de derde graad ($N=600$); steekproef: de 30 bevraagde leerlingen ($n=30$); variabele: slaapduur op een schoolnacht (uur); kwantitatief, continu. b) $\overline{x}\approx 7{,}19$ uur, $s\approx 0{,}94$ uur, $\text{Med}=7{,}25$ uur, $Q_1=6{,}7$ uur, $Q_3=7{,}7$ uur. c) Het histogram is bij benadering klokvormig en symmetrisch, met één top rond 7 – 7,5 uur; $\overline{x}\approx Med$; 21 van de 30 waarden (70 \%) liggen binnen $\overline{x}\pm s$, dat past bij de 68 \%-regel. Model: $N(7{,}19;\,0{,}94)$. d) $P(X<6)\approx 0{,}1024$; verwacht $600\cdot 0{,}1024\approx 61$ leerlingen. Antwoordzin: ongeveer 10 \% van de leerlingen, dus zowat 61 van de 600, slaapt minder dan 6 uur. e) Nee, het is een gemakssteekproef (enkel wie in de leefruimte zit; sociaalbewuste of net niet-slaperige leerlingen kunnen oververtegenwoordigd zijn). Beter: loten uit de volledige leerlingenlijst. f) Samenhang is geen causaliteit. Alternatieven: een verborgen variabele (bv. motivatie, thuissituatie, gezondheid) of omgekeerde causaliteit (wie goede punten heeft, heeft minder stress en slaapt beter). Onderzoek: experiment met testgroep en controlegroep, toevallig ingedeeld. g) Open antwoord: let op correcte cijfers (7,2 uur gemiddeld, standaardafwijking 0,9 uur, ongeveer 10 \% slaapt minder dan 6 uur), de vermelding dat de steekproef vertekend kan zijn en een concreet advies (bv. een nieuw onderzoek met een gelote steekproef).
\end{oplossing}

\subsection*{Eindmissie 2 · Twee korte problemen}

\subsubsection*{Probleem A \quad ★★★ \quad Een klasgemiddelde}

In een klas van 20 leerlingen is het gemiddelde op een toets 12,5 op 20. Twee leerlingen met 18 en 17 verlaten de klas. Wat is nu het gemiddelde van de 18 overblijvende leerlingen? Werk met de heuristiek van 7.1. \lijn

\begin{oplossing}
\textbf{Oplossing Probleem A}

Totaal: $20\cdot 12{,}5=250$. Zonder de twee leerlingen: $250-18-17=215$. Nieuw gemiddelde: $\dfrac{215}{18}\approx 11{,}9$ op 20. Antwoordzin: het gemiddelde van de 18 overblijvende leerlingen is ongeveer 11,9 op 20; het daalt van 12,5 omdat twee leerlingen met een resultaat boven het gemiddelde vertrokken.
\end{oplossing}

\subsubsection*{Probleem B \quad ★★★ \quad Een speeltuin}

Voor een attractie in een speeltuin moeten kinderen minstens 120 cm en hoogstens 140 cm groot zijn. De lengte van 8-jarigen is normaal verdeeld met $\mu=130$ cm en $\sigma=6$ cm.

\textbf{a)} Welk deel van de 8-jarigen mag de attractie gebruiken? \lijn

\textbf{b)} Er komt een groep van 150 kinderen van 8 jaar. Hoeveel kinderen verwacht je die te klein zijn voor de attractie? En hoeveel die te groot zijn? \lijn

\textbf{c)} Controleer je antwoord met de 68 – 95 – 99,7-regel. \lijn

\begin{oplossing}
\textbf{Oplossing Probleem B}

a) $P(120<X<140)\approx 0{,}9044$: ongeveer 90,4 \% van de 8-jarigen mag de attractie gebruiken. b) Te klein: $P(X<120)\approx 0{,}0478$, dus $150\cdot 0{,}0478\approx 7$ kinderen. Te groot: $P(X>140)\approx 0{,}0478$ (symmetrie), dus ook ongeveer 7 kinderen. c) De grenzen 120 en 140 liggen op $\mu\pm 1{,}67\sigma$, dus tussen $\mu\pm\sigma$ (68 \%) en $\mu\pm 2\sigma$ (95 \%). Het antwoord 90,4 \% ligt tussen 68 \% en 95 \%, dus is realistisch.
\end{oplossing}

\subsection*{Eindmissie 3 · Uitleggen en verbeteren (communiceren)}

\textbf{1.} Leg in je eigen woorden aan een klasgenoot uit waarom een gemiddelde alleen niet voldoende is om twee groepen te vergelijken. Gebruik een voorbeeld. \lijn

\textbf{2.} Een krant schrijft: ‘Uit een online poll met 8\,000 deelnemers blijkt dat 90 \% van de jongeren dagelijks een uur sport.’ Schrijf een korte, onderbouwde reactie (drie zinnen) met de begrippen uit dit boekje. \lijn

\textbf{3.} Bij een kansvraag schrijft een leerling enkel ‘0,1587’ als antwoord. Wat ontbreekt er volgens het stappenplan van level 6? \lijn

\begin{oplossing}
\textbf{Oplossing Eindmissie 3}

1. Voorbeeld: An en Bram hebben hetzelfde gemiddelde (25 min per bezoek), maar de standaardafwijking verschilt sterk (3,7 min tegenover 11,4 min); het gemiddelde zegt niets over hoe sterk de bezoeken variëren. 2. Online poll = zelfselectie, dus een vertekende steekproef; het grote aantal deelnemers maakt de steekproef niet representatief; enkel sportieve jongeren reageren waarschijnlijk. Het besluit mag dus niet veralgemeend worden naar alle jongeren. 3. Er ontbreken: het model $X\sim N(\mu,\sigma)$ met de context, een schets met de gearceerde oppervlakte, de ICT-berekening (welke functie en welke grenzen) en een antwoordzin die zegt wat de kans betekent. Ook een controle of het antwoord realistisch is.
\end{oplossing}
